% =============================================================================
% Interim report (Introduction, Literature Review, Problem Statement), IEEE two-column.
% SINGLE-FILE version for Overleaf, exported on 2026-11-15 from the repository
% (paper/interim.tex + paper/sections/*.tex) by scripts/export_overleaf.py.
%
% HOW TO USE IN OVERLEAF
%   1. New Project -> Blank Project, then replace the whole content of main.tex with this file.
%   2. Menu -> Compiler: pdfLaTeX (the default), then press Recompile. No other file is needed:
%      IEEEtran is built into Overleaf, all numbers are written in the text, and the
%      references are in the thebibliography block at the end.
%   3. To change a reference, edit its \bibitem entry at the end; to add a citation, add a
%      \bibitem with a new key and use \cite{key} in the text.
%
% The text was drafted with an AI assistant and has to be revised by the authors in their own
% words (see the Acknowledgment). Delete these comment lines before submitting.
% =============================================================================
\documentclass[10pt,conference]{IEEEtran}
\IEEEoverridecommandlockouts
\usepackage[utf8]{inputenc}
\usepackage[T1]{fontenc}
\usepackage{cite}
\usepackage{amsmath,amssymb}
\usepackage{array}
\usepackage{booktabs}
\usepackage[hidelinks]{hyperref}

\begin{document}

\title{Noise Drift of Simulator-Trained Quantum Classifiers on IBM Heron Processors:
Introduction, Literature Review and Problem Statement}

\author{
\IEEEauthorblockN{Pankaj Baid, Kartikey Gupta and Shreshtha Gupta}
\IEEEauthorblockA{Newton School of Technology, Rishihood University\\
Sonipat, Haryana, India\\
Corresponding author: pankaj.baid2024@nst.rishihood.edu.in}}

\maketitle

\noindent\textit{Interim report. It contains the Introduction, the Literature Review and the
Problem Statement of the full paper. The abstract, methodology, results and discussion, and
conclusion are in the full paper.}

\section{Introduction}
\label{sec:intro}

Most quantum machine learning (QML) models are designed and trained in a noiseless simulator
and only later evaluated on a quantum processor. On current hardware this hand-over is not
neutral. Gate errors, readout errors and decoherence change the expectation values that a
classifier reads out \cite{wang2021nibp,heyraud2022noisykernels}, and the noise is not fixed
either: cloud processors are recalibrated regularly and their parameters move between
calibrations \cite{realtimecal2023}. A model that works on the day it is tested may behave
differently a few days later, or on a sister device.

Earlier work has measured and narrowed the simulator-to-hardware gap at one point in time.
QuantumNAT showed that noise acts on measured outputs roughly as a scale and a shift, and
corrected for it with normalisation and noise-aware training \cite{wang2022quantumnat}.
Readout-error mitigation targets the bias that measurement errors add to expectation values
\cite{nation2021m3,vandenberg2022trex}. For uncertainty, conformal prediction turns a few
measurement shots into prediction sets with a coverage guarantee \cite{park2023qcp}, but noise
that changes over time can void the guarantee; an adaptive variant restores average coverage,
shown with data taken on one processor on a single day \cite{spencer2026aqcp}. Over longer
periods, a study of QAOA on three Heron processors found that emulators built from calibration
data overstated the hardware benefit of calibration-aware compilation \cite{qaoa2026drift}, and
a hardware benchmark on biosignal tasks found that test sets of the size a small budget allows
give unstable estimates \cite{qbiobench2026}. Section~\ref{sec:litreview} reviews this work.

For classifiers, a measurement is still missing, to our knowledge, that repeats the same frozen
models over several days and on more than one device, separates day-to-day variation from shot
noise, and prices each repair in QPU time. The last point matters because free cloud access,
such as IBM's Open Plan, provides ten minutes of processor time per 28
days\footnote{IBM Quantum plans:
\url{https://quantum.cloud.ibm.com/docs/en/guides/plans-overview}.}; any repair that needs
extra circuits competes with the experiment itself.

This paper reports such a measurement. We train three small classifiers on two datasets in
simulation, freeze them, and run the same 182 inputs per dataset on two IBM
Heron r2 processors, \texttt{ibm\_fez} and \texttt{ibm\_kingston}, in 5 sessions per device over
16 days. Packing 14 copies of each 4-qubit circuit onto
disjoint groups of qubits keeps a session near 20.5~s of QPU time, so the
whole campaign fits one free account. Hypotheses, analysis and decision rules were registered
before the first campaign session. We ask four questions:
\begin{itemize}
\item[RQ1] How large is the accuracy gap between simulator and hardware, and how well does an
emulator built from the reported calibration predict it?
\item[RQ2] How much do results vary between sessions and devices beyond shot noise, and which
reported calibration quantities explain the variation?
\item[RQ3] Does the scale-and-shift noise model hold on Heron processors, and is it stable
over time?
\item[RQ4] Which low-cost repairs recover accuracy and calibrated uncertainty, per QPU second?
\end{itemize}

Our contributions are:
\begin{itemize}
\item a pre-registered, budget-capped protocol for tracking QML classifiers on cloud hardware,
with circuit packing checked for crosstalk in a pilot;
\item paired measurements on two devices that answer RQ1--RQ4;
\item a comparison of four repair families by the accuracy they gain per QPU second, and of
four ways to calibrate conformal prediction sets;
\item the code, frozen models and raw data, with calibration snapshots taken at submission and
at job start.
\end{itemize}

We make no claim of quantum advantage. In simulation, tuned classical models match or beat the
quantum models: on two curves an RBF support vector machine reaches
0.987 against 0.927 for the best quantum model,
and on breast cancer the best accuracies differ by less than their confidence
intervals, in line with recent benchmark studies
\cite{bowles2024benchmarking,schuld2022advantage}. Our question is how reliably a given quantum
model behaves once it leaves the simulator. In brief, hardware accuracy was 9.2
points below the simulator on average,
the variance of the features between sessions was 3.8 times the variance expected from
shot noise,
and two of the eight repairs we tested gave a clear gain while the others did not.

The remainder of this report reviews the literature
(Section~\ref{sec:litreview}) and states the problem (Section~\ref{sec:problem}); the
methodology, the results and the conclusion are part of the full paper.

\section{Literature Review}
\label{sec:litreview}

This section reviews the work on which the study builds. It describes the quantum classifiers
that are run on hardware, how noise affects quantum models, the methods proposed to adapt to
or correct noise, ways to quantify uncertainty, and critical benchmarks. It ends with a
comparison of the closest studies and the gap that remains.

\subsection{Quantum classifiers on near-term hardware}
Early hardware demonstrations of quantum classification used a superconducting processor to
implement a variational classifier and a kernel estimator based on a quantum feature map
\cite{havlicek2019supervised}. Later work clarified what such models can represent. A model
that re-uploads its input several times can act as a universal classifier even on a single
qubit \cite{perezsalinas2020reuploading}. More generally, a circuit with data-encoding gates is
a partial Fourier series in the input, whose accessible frequencies are set by the encoding
\cite{schuld2021encoding}, and the standard circuit-based models, including re-uploading
circuits, can be written as linear models in a quantum feature space
\cite{jerbi2023beyondkernels}. Projecting the quantum state back to a classical representation
gives ``projected'' models, and the same work shows that, with training data available,
classical learners can stay competitive even on problems tailored to quantum computers
\cite{huang2021power}. A further family uses a fixed quantum circuit as a feature generator and
trains only a classical readout, as in quantum reservoir computing, which has been
demonstrated on superconducting devices \cite{yasuda2023qrc}. Our three classifiers belong to
these families: a re-uploading circuit (M1), projected single-qubit features (M2) and a fixed
random circuit with a linear readout (M3).

\subsection{Noise and the failure of quantum models}
Noise harms quantum models in several ways. For variational circuits, local Pauli noise makes
the gradient vanish exponentially with circuit depth, a noise-induced barren plateau
\cite{wang2021nibp}. Kernel methods face a related problem: kernel values can concentrate
exponentially around a constant, and noise is one of four identified sources, next to the
expressivity of the embedding, global measurements and entanglement
\cite{thanasilp2024concentration}. Decoherence also limits the number of independent features
that a noisy quantum kernel can extract, and acts as an implicit regulariser
\cite{heyraud2022noisykernels}. Remedies studied for kernels include tuning the bandwidth
\cite{shaydulin2022bandwidth}, although tuned quantum kernels can then resemble classical
radial-basis kernels \cite{bandwidth2025similarity}; training the embedding while mitigating
the effect of device noise \cite{hubregtsen2022embedding}; covariant kernels with a bit-flip
tolerance \cite{covariant2025concentration}; local and multi-scale kernels
\cite{localkernels2026}; and allocating a limited shot budget to the kernel entries that matter
most \cite{aqka2026}. These works concern training, kernel estimation and scaling with the
number of qubits. Our setting is different: the models are small, they are trained only in
simulation, and the question is how their predictions change when they are executed on hardware.

\subsection{Adapting models to device noise}
Several methods make a model robust to the noise of a specific device. QuantumNAT reported a gap
of more than 60\% between noise-free simulation and hardware for a four-class digit task on
IBMQ-Yorktown, found that noise acts on the measured outputs roughly as a scaling and a shift,
and countered it with post-measurement normalisation, noise injection during training and
quantisation of the outputs, improving accuracy by up to 43\% on six devices
\cite{wang2022quantumnat}. QuantumNAS searches jointly for a noise-robust circuit and qubit
mapping, and its authors note that a searched circuit must be evaluated on the machine right
after the search, because otherwise the noise characteristics drift \cite{wang2022quantumnas}.
QOC trains the parameters on the device with the parameter-shift rule and prunes the gradients
that noise has made unreliable \cite{wang2022qoc}. Across devices, a recent study transfers a
learned noise model from one IBM processor to another and finds that two-qubit gate error is
the main source of mismatch \cite{farib2026crossdevice}. Because superconducting processors
fluctuate, calibration has to be repeated, and real-time calibration frameworks aim to do so
without taking the processor offline \cite{realtimecal2023}.

\subsection{Error mitigation}
Error mitigation corrects expectation values after the fact, usually at the price of extra
circuits or shots. Readout errors, which bias every measured expectation value, can be mitigated
in a reduced subspace of the observed bit strings \cite{nation2021m3} or without a noise model,
by turning the bias into a multiplicative factor that is measured and removed
\cite{vandenberg2022trex}. For gate errors, zero-noise extrapolation and probabilistic error
cancellation are standard techniques for short circuits \cite{temme2017zne}; Clifford data
regression learns the correction from near-Clifford circuits that can be simulated classically
\cite{czarnik2021cdr}; and machine-learned mitigation lowers the cost of mitigation, with its
authors listing the drift of noise in hardware as an open question \cite{liao2024mlqem}. On a
free budget the extra runtime matters, because mitigation circuits compete with the experiment
for the same processor time. For this reason we test only repairs that need no additional
circuits, a small number of calibration circuits, or a few labelled hardware samples.

\subsection{Uncertainty and calibration of quantum predictions}
A classifier should also report how far its predictions can be trusted. Calibration measures such
as the expected calibration error were popularised for neural networks \cite{guo2017calibration},
and a recent library provides shot-based and ensemble-based uncertainty estimates for quantum
models and reports that variational classifiers tend to be overconfident \cite{quantumuq2026}.
Finite-shot evaluation by itself randomises a deployed model, and an audit of frozen quantum
kernel classifiers found that this moves thresholded metrics by about 0.02
\cite{fernandezbarrios2026conditional}. Conformal prediction turns a few shots of a quantum
model into prediction sets with a coverage guarantee \cite{park2023qcp}. The guarantee assumes
exchangeable data, which time-varying hardware noise can violate. Adaptive conformal inference
\cite{gibbs2021aci} lifts the assumption by updating the working miscoverage level online, and
AQCP applies it to quantum models, with hardware data from a single day and a multi-day
assessment left to future work \cite{spencer2026aqcp}.

\subsection{Benchmarking and the evidence for quantum advantage}
Claims that quantum models outperform classical ones need careful baselines. A large study of
12 popular models on 6 binary classification tasks, with 160 datasets, found that
out-of-the-box classical models outperform the quantum classifiers and that removing
entanglement often leaves performance unchanged or better \cite{bowles2024benchmarking}. The
same study counts that about 40\% of a sample of 55 arXiv papers claim that a quantum model
beats a classical one, whereas only 3 report the opposite. A related argument holds that a focus
on advantage restricts attention to a biased subset of models and datasets
\cite{schuld2022advantage}. Quantum convolutional networks look strong on phase-recognition
benchmarks, yet on such locally easy datasets they can be simulated classically
\cite{bermejo2026qcnn}, and an attribution audit of intrusion-detection benchmarks traced
apparent gains to classical preprocessing and regularisation \cite{gillani2026nids}. On
hardware, QBioBench evaluated kernel and variational classifiers on an IBM Nighthawk processor:
no quantum model reached the performance of its matched support vector machine, readout
mitigation changed the ROC-AUC by at most 0.015, and random subsets of the test pool showed
that the small test sets favoured by circuit budgets give unreliable estimates
\cite{qbiobench2026}. We take two lessons from this literature: report tuned classical
baselines, and use a test set that is fixed across sessions.

\subsection{Hardware studies of calibration, drift and parallel execution}
Few studies follow the same workload on hardware over days. The closest is a study of QAOA for
weighted MaxCut on three IBM Heron r2 processors, which combined a 54-pair hardware grid,
repeats of the same circuits 9.3 to 12.1 days later and a refresh experiment. Selecting qubits
with the full calibration improved the median approximation ratio by 0.74 percentage points; an
emulator built from the calibration reproduced the direction of this advantage in 35 of 54
comparisons but overstated its size, and the compiler effect kept its sign in 16 of 18 repeats
\cite{qaoa2026drift}. The authors conclude that calibration data are useful for finding
candidates and deciding when to refresh, but that scalar calibration scores do not predict
hardware performance well. Running many small circuits side by side raises throughput but can
add crosstalk, which has to be characterised and scheduled around \cite{murali2020crosstalk}.
The palloq method chooses which circuits to combine and where to place them to limit this
interference \cite{ohkura2022simultaneous}, and a study of parallel QML circuits on an IBM
processor measures how the occupancy of the processor affects the loss of the learning
algorithm \cite{parallelqml2024}. Our packing of 14 circuit copies per job follows
this line of work, with a pilot run to check for crosstalk. All circuits are built and
transpiled with Qiskit \cite{javadiabhari2024qiskit}.

\begin{table*}[t]
\centering
\caption{Closest prior studies and what remains open. Statements about other studies follow
their abstracts or text.}
\label{tab:related}
\footnotesize
\setlength{\tabcolsep}{4pt}
\renewcommand{\arraystretch}{1.15}
\begin{tabular}{>{\raggedright\arraybackslash}p{2.5cm}>{\raggedright\arraybackslash}p{3.3cm}>{\raggedright\arraybackslash}p{3.3cm}>{\raggedright\arraybackslash}p{3.2cm}>{\raggedright\arraybackslash}p{3.9cm}}
\toprule
Study & Workload & Hardware and period & Correction or uncertainty method & What remains open \\
\midrule
QuantumNAT \cite{wang2022quantumnat} & Quantum neural networks, 8 classification tasks & 6 IBM devices & Normalisation, noise injection, quantisation & Variation across days is not reported \\
QuantumNAS \cite{wang2022quantumnas} & Joint search of circuit and qubit mapping & IBM devices & Search under a noise model & Noise drifts, so a searched circuit must be run right after the search \\
AQCP \cite{spencer2026aqcp} & Probabilistic regression, 5 qubits & \texttt{ibm\_sherbrooke}, one day & Adaptive conformal inference & Multi-day assessment left to future work \\
QAOA drift study \cite{qaoa2026drift} & Weighted MaxCut with QAOA & 3 Heron r2 processors, repeats after 9--12 days & Calibration-aware qubit selection and refresh & Optimisation, not classification; emulator overstated the gain \\
QBioBench \cite{qbiobench2026} & Biosignal classification, kernels and variational circuits & One IBM Nighthawk r2 processor & Readout mitigation & Single processor, exploratory; small test sets unreliable \\
Cross-device transfer \cite{farib2026crossdevice} & Learning noise models of circuit outputs & \texttt{ibm\_fez} and \texttt{ibm\_marrakesh} & Few-shot fine-tuning, 20 samples & Targets output distributions, not classifier accuracy \\
Learned mitigation \cite{liao2024mlqem} & Expectation values, up to 100 qubits & IBM devices & Machine-learned mitigation & Drift of noise named as an open question \\
\midrule
\textbf{This study} & Three classifiers, two datasets & Two Heron r2 processors, 10 sessions over 16 days & Four repair families and four conformal schemes, each priced in QPU seconds & Frozen small models; free-plan budget \\
\bottomrule
\end{tabular}
\end{table*}

\subsection{Summary and research gap}
\label{sec:gap}
Table~\ref{tab:related} compares the closest studies. Four points emerge. First, evidence on how
deployed classifiers change over days is thin: QuantumNAS and the learned-mitigation study both
mention that noise drifts \cite{wang2022quantumnas,liao2024mlqem}, yet the multi-day hardware
study we found uses QAOA \cite{qaoa2026drift}, and the hardware evidence for conformal
prediction under time-varying noise comes from one day \cite{spencer2026aqcp}. Second,
corrections are mostly evaluated on one device at a time; the cross-device work we found
concerns noise models rather than classifier accuracy \cite{farib2026crossdevice}. Third, the
cost of a correction is usually described qualitatively, although on a free plan it decides what
can be run at all. Fourth, small hardware test sets make the accuracy of a single session
unreliable \cite{qbiobench2026}, so effects have to be pooled over a fixed test set and
repeated sessions.

To our knowledge, based on searches of OpenAlex, Semantic Scholar and arXiv in October and
November 2026, no earlier study repeats frozen quantum classifiers on more than one IBM Heron
processor over several days, separates drift from shot noise, compares several corrections by
their QPU cost, and checks conformal coverage on the same data. This is the gap that the
present work addresses.

\section{Problem Statement}
\label{sec:problem}

\textbf{Problem.} A quantum classifier is trained in a noiseless simulator and then frozen. How
much do its accuracy and the reliability of its predictions change when the same circuits run
on real IBM Heron processors on different days and devices? How much of the change is genuine
drift rather than shot noise? And which low-cost corrections recover the loss for the least QPU
time? The gap in the literature that motivates these questions is summarised in
Section~\ref{sec:gap}.

\subsection{Setting}
Let $x\in[0,\pi]^{n}$ be an input with $n=4$ features and $\rho_\theta(x)$ the state
that a fixed circuit prepares on $n$ qubits. A model reads out Pauli expectation values and
passes them to a logistic head,
\begin{align}
\hat p(y=1\mid x)&=\sigma\!\Big(w^\top\frac{\phi(x)-\mu}{s}+b\Big), \label{eq:model}\\
\phi_j(x)&=\operatorname{tr}\!\big[P_j\,\rho_\theta(x)\big], \nonumber
\end{align}
where $P_j$ is a Pauli string and $\mu$, $s$ are the mean and standard deviation of the training
features. The circuit parameters $\theta$ and the head $(w,b,\mu,s)$ are fixed after training in
simulation. On device $d$ in session $t$, a job runs the circuits with $N=500$ shots and
returns estimates $\hat\phi_{d,t}(x)$. Their expectation $\bar\phi_{d,t}(x)$ differs from
$\phi(x)$ because of gate, readout and decoherence errors, and it changes with $t$ as the
processor drifts and is recalibrated. Shot noise adds a variance of
$(1-\bar\phi_{d,t,j}^{2})/N$ to feature $j$. Accuracy is measured on a fixed test set of
150 inputs, giving $A_{\mathrm{sim}}$ in the simulator and $A_{d,t}$ on hardware.

\subsection{Quantities to determine}
\begin{itemize}
\item \emph{Deployment gap and emulator error.} $G_{d,t}=A_{\mathrm{sim}}-A_{d,t}$ and
$E_{d,t}=A^{\mathrm{emu}}_{d,t}-A_{d,t}$, where $A^{\mathrm{emu}}_{d,t}$ is the accuracy
predicted by a noise model built from the calibration reported at submission.
\item \emph{Drift.} For each input and feature, the ratio of the variance of the estimate across
sessions to its expected shot variance,
\begin{equation}
r_d(x,j)=\frac{\operatorname{Var}_t\big[\hat\phi_{d,t,j}(x)\big]}
{\operatorname{E}_t\big[(1-\hat\phi_{d,t,j}(x)^{2})/N\big]}.
\label{eq:drift}
\end{equation}
A ratio of one means that shot noise alone explains the variation between sessions.
\item \emph{Noise map.} The scale $\gamma$ and shift $\beta$ in
$\bar\phi_{d,t,j}(x)\approx\gamma\,\phi_j(x)+\beta$, and their stability over $t$.
\item \emph{Repairs.} A repair $\mathcal R$ maps the hardware counts, and possibly $k$ labelled
hardware inputs, to corrected predictions. Its gain is $g(\mathcal R)=A^{\mathcal R}_{d,t}-
A_{d,t}$, and its cost $c(\mathcal R)$ is the extra QPU seconds per session.
\item \emph{Uncertainty.} Prediction sets $C_\alpha(x)\subseteq\{0,1\}$ with target coverage
$1-\alpha=0.90$, judged by the empirical coverage
$\Pr[y\in C_\alpha(x)]$ and the mean set size.
\end{itemize}

\subsection{Research questions and hypotheses}
Table~\ref{tab:rq} links each research question to the quantity above and to the hypothesis that
was registered after the pilot runs and before the first campaign session (git tag
\texttt{prereg-v1}, 2026-10-27). Every hypothesis has a decision rule, fixed in advance and
based on 95\% bootstrap intervals, so that its outcome does not depend on the data.

\begin{table}[t]
\centering
\caption{Research questions, the quantities measured and the registered hypotheses.}
\label{tab:rq}
\footnotesize
\setlength{\tabcolsep}{3pt}
\begin{tabular}{@{}l>{\raggedright\arraybackslash}p{2.35cm}>{\raggedright\arraybackslash}p{4.2cm}@{}}
\toprule
 & Quantity & Registered hypothesis \\
\midrule
RQ1 & $G_{d,t}$, $E_{d,t}$ & H1: the emulator overestimates hardware accuracy ($E_{d,t}>0$) \\
RQ2 & drift ratio $r_d$; regression on chain calibration & H2: the ratio exceeds one. H4: readout error is the strongest predictor of feature error \\
RQ3 & $\gamma$, $\beta$ per session & none (descriptive) \\
RQ4 & $g(\mathcal R)$, $c(\mathcal R)$, coverage & H3: a 16-label refit beats normalisation. H5: split conformal calibrated in session 1 under-covers later sessions, while ACI holds \\
\bottomrule
\end{tabular}
\end{table}

\subsection{Constraints and scope}
The experiments use the free IBM Open Plan, which gives ten minutes of processor time per 28
days. The campaign is capped at 400~s of QPU time, jobs run when a shared queue
allows, and the models, datasets and qubit chains are fixed before the first session. We study
small classifiers on classical data (4 qubits, two datasets) and do not train on
hardware. We do not claim quantum advantage: the quantum models are compared with tuned
classical baselines only to put their accuracy in context. The study is a measurement of
deployment behaviour on two processors of one family, not a general model of quantum noise.

\section*{Acknowledgment}
We acknowledge the use of IBM Quantum services for this work. The views expressed are those of
the authors, and do not reflect the official policy or position of IBM or the IBM Quantum team.
The analysis software and first drafts of parts of this manuscript were prepared with the help
of an AI assistant (Kiro); the authors reviewed, revised and verified all code, text and results.

% Generated by IEEEtran.bst, version: 1.14 (2015/08/26)
\begin{thebibliography}{10}
\providecommand{\url}[1]{#1}
\csname url@samestyle\endcsname
\providecommand{\newblock}{\relax}
\providecommand{\bibinfo}[2]{#2}
\providecommand{\BIBentrySTDinterwordspacing}{\spaceskip=0pt\relax}
\providecommand{\BIBentryALTinterwordstretchfactor}{4}
\providecommand{\BIBentryALTinterwordspacing}{\spaceskip=\fontdimen2\font plus
\BIBentryALTinterwordstretchfactor\fontdimen3\font minus
  \fontdimen4\font\relax}
\providecommand{\BIBforeignlanguage}[2]{{%
\expandafter\ifx\csname l@#1\endcsname\relax
\typeout{** WARNING: IEEEtran.bst: No hyphenation pattern has been}%
\typeout{** loaded for the language `#1'. Using the pattern for}%
\typeout{** the default language instead.}%
\else
\language=\csname l@#1\endcsname
\fi
#2}}
\providecommand{\BIBdecl}{\relax}
\BIBdecl

\bibitem{wang2021nibp}
S.~Wang, E.~Fontana, M.~Cerezo, K.~Sharma, A.~Sone, L.~Cincio, and P.~J. Coles,
  ``{Noise-induced barren plateaus in variational quantum algorithms},''
  \emph{Nature Communications}, vol.~12, p. 6961, 2021, doi:
  \href{https://doi.org/10.1038/s41467-021-27045-6}{10.1038/s41467-021-27045-6}.

\bibitem{heyraud2022noisykernels}
V.~Heyraud, Z.~Li, Z.~Denis, A.~Le~Boit{\'e}, and C.~Ciuti, ``{Noisy quantum
  kernel machines},'' \emph{Physical Review A}, vol. 106, p. 052421, 2022, doi:
  \href{https://doi.org/10.1103/PhysRevA.106.052421}{10.1103/PhysRevA.106.052421}.

\bibitem{realtimecal2023}
Y.~Liu, Z.~Li, A.~Robertson, X.~Fu, and S.~L. Song, ``{Enabling Efficient
  Real-Time Calibration on Cloud Quantum Machines},'' \emph{IEEE Transactions
  on Quantum Engineering}, vol.~4, pp. 1--17, 2023, doi:
  \href{https://doi.org/10.1109/TQE.2023.3276970}{10.1109/TQE.2023.3276970}.

\bibitem{wang2022quantumnat}
H.~Wang, J.~Gu, Y.~Ding, Z.~Li, F.~T. Chong, D.~Z. Pan, and S.~Han,
  ``{QuantumNAT: Quantum Noise-Aware Training with Noise Injection,
  Quantization and Normalization},'' in \emph{Proceedings of the 59th ACM/IEEE
  Design Automation Conference (DAC)}, 2022, doi:
  \href{https://doi.org/10.1145/3489517.3530400}{10.1145/3489517.3530400}.

\bibitem{nation2021m3}
P.~D. Nation, H.~Kang, N.~Sundaresan, and J.~M. Gambetta, ``{Scalable
  Mitigation of Measurement Errors on Quantum Computers},'' \emph{PRX Quantum},
  vol.~2, p. 040326, 2021, doi:
  \href{https://doi.org/10.1103/PRXQuantum.2.040326}{10.1103/PRXQuantum.2.040326}.

\bibitem{vandenberg2022trex}
E.~van~den Berg, Z.~K. Minev, and K.~Temme, ``{Model-free readout-error
  mitigation for quantum expectation values},'' \emph{Physical Review A}, vol.
  105, p. 032620, 2022, doi:
  \href{https://doi.org/10.1103/PhysRevA.105.032620}{10.1103/PhysRevA.105.032620}.

\bibitem{park2023qcp}
S.~Park and O.~Simeone, ``{Quantum Conformal Prediction for Reliable
  Uncertainty Quantification in Quantum Machine Learning},'' \emph{IEEE
  Transactions on Quantum Engineering}, 2023, doi:
  \href{https://doi.org/10.1109/TQE.2023.3333224}{10.1109/TQE.2023.3333224}.

\bibitem{spencer2026aqcp}
D.~Spencer, S.~Nicholls, and M.~Caprio, ``{Adaptive Conformal Prediction for
  Quantum Machine Learning},'' \emph{Transactions on Machine Learning
  Research}, 2026, arXiv:2511.18225.

\bibitem{qaoa2026drift}
A.~Vishwakarma, ``{When Does Calibration-Aware Compilation Transfer to Quantum
  Hardware? A Three-Processor QAOA Study of Ranking, Drift, and Refresh},''
  \emph{Quantum Science and Technology}, 2026, doi:
  \href{https://doi.org/10.1088/2058-9565/aeb8eb}{10.1088/2058-9565/aeb8eb}.

\bibitem{qbiobench2026}
M.~Sokol, J.~Hejda, P.~Volf, and P.~Kut{\'i}lek, ``{QBioBench: Benchmark of
  quantum machine learning for peripheral biosignals on real hardware},''
  \emph{Quantum Science and Technology}, 2026, doi:
  \href{https://doi.org/10.1088/2058-9565/aeb71e}{10.1088/2058-9565/aeb71e}.

\bibitem{bowles2024benchmarking}
J.~Bowles, S.~Ahmed, and M.~Schuld, ``{Better than classical? The subtle art of
  benchmarking quantum machine learning models},'' arXiv:2403.07059, 2024.

\bibitem{schuld2022advantage}
M.~Schuld and N.~Killoran, ``{Is Quantum Advantage the Right Goal for Quantum
  Machine Learning?}'' \emph{PRX Quantum}, vol.~3, p. 030101, 2022, doi:
  \href{https://doi.org/10.1103/PRXQuantum.3.030101}{10.1103/PRXQuantum.3.030101}.

\bibitem{havlicek2019supervised}
V.~Havl{\'i}{\v{c}}ek, A.~D. C{\'o}rcoles, K.~Temme, A.~W. Harrow, A.~Kandala,
  J.~M. Chow, and J.~M. Gambetta, ``{Supervised learning with quantum-enhanced
  feature spaces},'' \emph{Nature}, vol. 567, pp. 209--212, 2019, doi:
  \href{https://doi.org/10.1038/s41586-019-0980-2}{10.1038/s41586-019-0980-2}.

\bibitem{perezsalinas2020reuploading}
A.~P{\'e}rez-Salinas, A.~Cervera-Lierta, E.~Gil-Fuster, and J.~I. Latorre,
  ``{Data re-uploading for a universal quantum classifier},'' \emph{Quantum},
  vol.~4, p. 226, 2020, doi:
  \href{https://doi.org/10.22331/q-2020-02-06-226}{10.22331/q-2020-02-06-226}.

\bibitem{schuld2021encoding}
M.~Schuld, R.~Sweke, and J.~J. Meyer, ``{Effect of data encoding on the
  expressive power of variational quantum-machine-learning models},''
  \emph{Physical Review A}, vol. 103, p. 032430, 2021, doi:
  \href{https://doi.org/10.1103/PhysRevA.103.032430}{10.1103/PhysRevA.103.032430}.

\bibitem{jerbi2023beyondkernels}
S.~Jerbi, L.~J. Fiderer, H.~Poulsen~Nautrup, J.~M. K{\"u}bler, H.~J. Briegel,
  and V.~Dunjko, ``{Quantum machine learning beyond kernel methods},''
  \emph{Nature Communications}, vol.~14, p. 517, 2023, doi:
  \href{https://doi.org/10.1038/s41467-023-36159-y}{10.1038/s41467-023-36159-y}.

\bibitem{huang2021power}
H.-Y. Huang, M.~Broughton, M.~Mohseni, R.~Babbush, S.~Boixo, H.~Neven, and
  J.~R. McClean, ``{Power of data in quantum machine learning},'' \emph{Nature
  Communications}, vol.~12, p. 2631, 2021, doi:
  \href{https://doi.org/10.1038/s41467-021-22539-9}{10.1038/s41467-021-22539-9}.

\bibitem{yasuda2023qrc}
T.~Yasuda, Y.~Suzuki, T.~Kubota, K.~Nakajima, Q.~Gao, W.~Zhang, S.~Shimono,
  A.~Matsuo, J.~Ananthanarayanan, H.~I. Nurdin, and N.~Yamamoto, ``{Quantum
  reservoir computing with repeated measurements on superconducting devices},''
  arXiv:2310.06706, 2023.

\bibitem{thanasilp2024concentration}
S.~Thanasilp, S.~Wang, M.~Cerezo, and Z.~Holmes, ``{Exponential concentration
  in quantum kernel methods},'' \emph{Nature Communications}, vol.~15, p. 5200,
  2024, doi:
  \href{https://doi.org/10.1038/s41467-024-49287-w}{10.1038/s41467-024-49287-w}.

\bibitem{shaydulin2022bandwidth}
R.~Shaydulin and S.~M. Wild, ``{Importance of kernel bandwidth in quantum
  machine learning},'' \emph{Physical Review A}, vol. 106, p. 042407, 2022,
  doi:
  \href{https://doi.org/10.1103/PhysRevA.106.042407}{10.1103/PhysRevA.106.042407}.

\bibitem{bandwidth2025similarity}
R.~Flórez-Ablan, M.~Roth, and J.~Schnabel, ``{On the similarity of
  bandwidth-tuned quantum kernels and classical kernels},'' \emph{Quantum
  Science and Technology}, 2025, doi:
  \href{https://doi.org/10.1088/2058-9565/ade7ad}{10.1088/2058-9565/ade7ad}.

\bibitem{hubregtsen2022embedding}
T.~Hubregtsen, D.~Wierichs, E.~Gil-Fuster, P.-J. H.~S. Derks, P.~K. Faehrmann,
  and J.~J. Meyer, ``{Training quantum embedding kernels on near-term quantum
  computers},'' \emph{Physical Review A}, vol. 106, p. 042431, 2022, doi:
  \href{https://doi.org/10.1103/PhysRevA.106.042431}{10.1103/PhysRevA.106.042431}.

\bibitem{covariant2025concentration}
G.~Agliardi, G.~Cortiana, A.~Dekusar, K.~Ghosh, N.~Mohseni, C.~O'Meara,
  V.~Valls, K.~Yogaraj, and S.~Zhuk, ``{Mitigating exponential concentration in
  covariant quantum kernels for subspace and real-world data},'' \emph{npj
  Quantum Information}, 2025, doi:
  \href{https://doi.org/10.1038/s41534-025-01154-2}{10.1038/s41534-025-01154-2}.

\bibitem{localkernels2026}
C.~Zendejas-Morales, D.~Saikia, and U.~Singh, ``{Local and Multi-Scale
  Strategies to Mitigate Exponential Concentration in Quantum Kernels},''
  arXiv:2602.16097, 2026.

\bibitem{aqka2026}
J.~Xu, C.~Li, D.~Zeng, J.~Paisley, and Q.~Zhao, ``{AQKA: Active Quantum Kernel
  Acquisition Under a Shot Budget},'' arXiv:2605.14672, 2026.

\bibitem{wang2022quantumnas}
H.~Wang, Y.~Ding, J.~Gu, Y.~Lin, D.~Z. Pan, F.~T. Chong, and S.~Han,
  ``{QuantumNAS: Noise-Adaptive Search for Robust Quantum Circuits},'' in
  \emph{IEEE International Symposium on High-Performance Computer Architecture
  (HPCA)}, 2022, arXiv:2107.10845.

\bibitem{wang2022qoc}
H.~Wang, Z.~Li, J.~Gu, Y.~Ding, D.~Z. Pan, and S.~Han, ``{QOC: Quantum On-Chip
  Training with Parameter Shift and Gradient Pruning},'' in \emph{Proceedings
  of the 59th ACM/IEEE Design Automation Conference (DAC)}, 2022,
  arXiv:2202.13239.

\bibitem{farib2026crossdevice}
S.~A. Farib, S.~R. Islam, and A.~R. Anik, ``{Few-Shot Cross-Device Transfer for
  Quantum Noise Modeling on Real Hardware},'' arXiv:2604.24397, 2026.

\bibitem{temme2017zne}
K.~Temme, S.~Bravyi, and J.~M. Gambetta, ``{Error Mitigation for Short-Depth
  Quantum Circuits},'' \emph{Physical Review Letters}, vol. 119, p. 180509,
  2017, doi:
  \href{https://doi.org/10.1103/PhysRevLett.119.180509}{10.1103/PhysRevLett.119.180509}.

\bibitem{czarnik2021cdr}
P.~Czarnik, A.~Arrasmith, P.~J. Coles, and L.~Cincio, ``{Error mitigation with
  Clifford quantum-circuit data},'' \emph{Quantum}, vol.~5, p. 592, 2021, doi:
  \href{https://doi.org/10.22331/q-2021-11-26-592}{10.22331/q-2021-11-26-592}.

\bibitem{liao2024mlqem}
H.~Liao, D.~S. Wang, I.~Sitdikov, C.~Salcedo, A.~Seif, and Z.~K. Minev,
  ``{Machine learning for practical quantum error mitigation},'' \emph{Nature
  Machine Intelligence}, 2024, doi:
  \href{https://doi.org/10.1038/s42256-024-00927-2}{10.1038/s42256-024-00927-2}.

\bibitem{guo2017calibration}
C.~Guo, G.~Pleiss, Y.~Sun, and K.~Q. Weinberger, ``{On Calibration of Modern
  Neural Networks},'' in \emph{Proceedings of the 34th International Conference
  on Machine Learning (ICML)}, 2017, arXiv:1706.04599.

\bibitem{quantumuq2026}
F.~{\"O}. {\c{C}}atak, ``{QuantumUQ: A Library for Uncertainty Quantification
  in Quantum Machine Learning},'' \emph{Nordic Machine Intelligence}, vol.~6,
  no.~1, pp. 19--26, 2026, doi:
  \href{https://doi.org/10.5617/nmi.13217}{10.5617/nmi.13217}.

\bibitem{fernandezbarrios2026conditional}
R.~Fernández-Barrios, I.~Pastor-López, A.~González-Santocildes, and P.~G.
  Bringas, ``{Conditional validity of quantum event classifiers under collider
  systematics and quantum estimation uncertainty},'' arXiv:2609.02781, 2026.

\bibitem{gibbs2021aci}
I.~Gibbs and E.~Cand{\`e}s, ``{Adaptive Conformal Inference Under Distribution
  Shift},'' in \emph{Advances in Neural Information Processing Systems
  (NeurIPS)}, 2021, arXiv:2106.00170.

\bibitem{bermejo2026qcnn}
P.~Bermejo, P.~Braccia, M.~S. Rudolph, Z.~Holmes, L.~Cincio, and M.~Cerezo,
  ``{Quantum Convolutional Neural Networks are (Effectively) Classically
  Simulable},'' \emph{PRX Quantum}, vol.~7, p. 020304, 2026, doi:
  \href{https://doi.org/10.1103/8qt9-72ts}{10.1103/8qt9-72ts}.

\bibitem{gillani2026nids}
S.~A. Gillani, M.~S.~A. Baig, S.~M. Shah, A.~Ali, and H.~Siddiqui, ``{How
  Quantum Is the Advantage? A Fair, Calibration- and Noise-Aware Benchmark and
  Attribution Audit of Quantum Machine Learning for Network Intrusion
  Detection},'' arXiv:2608.18155, 2026.

\bibitem{murali2020crosstalk}
P.~Murali, D.~C. McKay, M.~Martonosi, and A.~Javadi-Abhari, ``{Software
  Mitigation of Crosstalk on Noisy Intermediate-Scale Quantum Computers},'' in
  \emph{Proceedings of ASPLOS}, 2020, doi:
  \href{https://doi.org/10.1145/3373376.3378477}{10.1145/3373376.3378477}.

\bibitem{ohkura2022simultaneous}
Y.~Ohkura, T.~Satoh, and R.~Van~Meter, ``{Simultaneous Execution of Quantum
  Circuits on Current and Near-Future NISQ Systems},'' \emph{IEEE Transactions
  on Quantum Engineering}, vol.~3, 2022, doi:
  \href{https://doi.org/10.1109/TQE.2022.3164716}{10.1109/TQE.2022.3164716}.

\bibitem{parallelqml2024}
J.~Rauch, B.~Chichereau, S.~Vialle, P.~Carribault, and D.~Rontani,
  ``{Investigating parallel execution of quantum Machine Learning circuits on
  superconducting hardware},'' in \emph{IEEE International Conference on
  Quantum Computing and Engineering (QCE)}, 2024, pp. 199--204, doi:
  \href{https://doi.org/10.1109/QCE60285.2024.10278}{10.1109/QCE60285.2024.10278}.

\bibitem{javadiabhari2024qiskit}
A.~Javadi-Abhari \emph{et~al.}, ``{Quantum computing with Qiskit},''
  arXiv:2405.08810, 2024.

\end{thebibliography}

\end{document}
